%%
%% kleinbottle.tex
%%
%% Copyright (C) 2026 Benjamim Lima Neto
%%
%% This work may be distributed and/or modified under
%% the conditions of the LaTeX Project Public License,
%% either version 1.3c of this license or (at your option)
%% any later version.
%%
%% The latest version of this license is available at:
%% https://www.latex-project.org/lppl.txt
%%
%% This work has the LPPL maintenance status `maintained'.
%%
%% The Current Maintainer of this work is Benjamim Lima Neto.
%%
%% Contact information is provided in the accompanying
%% README.md file.
%%
%% This work consists of the files:
%%   kleinbottle.sty
%%   kleinbottle.tex
%%   README.md
%%
%% and the derived file:
%%   kleinbottle.pdf
%%

\documentclass[12pt,a4paper]{article}

% --------------------------------------------------
% Packages
% --------------------------------------------------

\usepackage{iftex}

\ifPDFTeX
    \usepackage[utf8]{inputenc}
    \usepackage[T1]{fontenc}
\fi

\usepackage{lmodern}
\usepackage[margin=2.5cm]{geometry}
\usepackage{array}
\usepackage{booktabs}
\usepackage{xcolor}
\usepackage{kleinbottle}
\usepackage{hyperref}

% --------------------------------------------------
% Colors
% --------------------------------------------------

\definecolor{ManualBlue}{HTML}{234A76}
\definecolor{CustomGreen}{HTML}{2A9D8F}
\definecolor{CustomPurple}{HTML}{7952B3}

% --------------------------------------------------
% Hyperlinks and PDF metadata
% --------------------------------------------------

\hypersetup{
    colorlinks=true,
    linkcolor=ManualBlue,
    urlcolor=ManualBlue,
    citecolor=ManualBlue,
    pdftitle={The kleinbottle package},
    pdfauthor={Benjamim Lima Neto},
    pdfsubject={Documentation of the kleinbottle LaTeX package},
    pdfkeywords={LaTeX, TikZ, Klein bottle, topology}
}

% --------------------------------------------------
% Title
% --------------------------------------------------

\title{%
    \leavevmode
    \klein[color=ManualBlue,scale=5]
    \par\vspace{0.4cm}
    \texttt{kleinbottle}
    \par\vspace{0.2cm}
    {\large A Tiny \LaTeX{} Package for Drawing a Klein Bottle}
}

\author{Benjamim Lima Neto}

\date{Version 1.0 -- November 26, 2026}

% --------------------------------------------------
% Document
% --------------------------------------------------

\begin{document}

\maketitle
\begin{abstract}

The \texttt{kleinbottle} package provides a simple and
customizable TikZ illustration of the Klein bottle.
Through a single command, users can insert a miniature
representation of this mathematical object into their
\LaTeX{} documents and customize its color, opacity,
size, and rotation.

The package is designed to be lightweight, easy to use,
and suitable for both inline and displayed illustrations.

\end{abstract}

\tableofcontents

% --------------------------------------------------
\section{Introduction}
% --------------------------------------------------

The Klein bottle is a well-known example of a compact
non-orientable surface without boundary. It is named
after the German mathematician Felix Klein.

The \texttt{kleinbottle} package provides a stylized
two-dimensional representation of the usual
self-intersecting immersion of the Klein bottle
in three-dimensional Euclidean space.

Its purpose is to make it easy to incorporate this
mathematical object into articles, presentations,
lecture notes, books, and other \LaTeX{} documents,
without requiring users to write their own TikZ code.

The package is intentionally lightweight and provides
a single command with four independent customization
options.

The illustration is generated directly from vector
graphics instructions and does not require external
image files.

% --------------------------------------------------
\section{Requirements}
% --------------------------------------------------

The package requires the following packages:

\begin{itemize}

    \item \href{https://ctan.org/pkg/pgf}
        {\texttt{TikZ/PGF}}:
        for drawing the Klein bottle, processing
        options, and performing numerical calculations;

    \item \href{https://ctan.org/pkg/xcolor}
        {\texttt{xcolor}}:
        for color customization;

    \item \href{https://ctan.org/pkg/graphicx}
        {\texttt{graphicx}}:
        for proportional scaling.

\end{itemize}

All dependencies are loaded automatically by
\texttt{kleinbottle}. There is no need to load them
manually unless additional functionality is required.

The package is intended for use with modern
\LaTeX{} installations, including those using
pdf\LaTeX, Xe\LaTeX, and Lua\LaTeX.

% --------------------------------------------------
\section{Installation}
% --------------------------------------------------

For manual installation, place the file
\texttt{kleinbottle.sty} in the same directory as your
main \LaTeX{} document, or in an appropriate local
TeX directory.

Then load the package in the preamble:

\begin{verbatim}
\usepackage{kleinbottle}
\end{verbatim}

No additional configuration is required.

If the package is already installed in your TeX
distribution, loading it in the preamble is sufficient.

% --------------------------------------------------
\section{Basic Usage}
% --------------------------------------------------

To insert the default Klein bottle miniature, simply
write:

\begin{verbatim}
\klein
\end{verbatim}

The result is:

\begin{center}
    \klein
\end{center}

The command accepts an optional argument containing
a comma-separated list of key--value pairs.

Its complete syntax is:

\begin{verbatim}
\klein[color=<color>,
       opacity=<number>,
       scale=<number>,
       rotate=<angle>]
\end{verbatim}

All four options are optional and may be specified
individually, simultaneously, or in any order.

Whenever an option is omitted, its default value
is used.

An empty optional argument is equivalent to omitting
the optional argument entirely:

\begin{verbatim}
\klein
\klein[]
\end{verbatim}

\begin{center}
    \klein
    \qquad
    \klein[]
\end{center}

\subsection{Available Options}

\begin{center}
\renewcommand{\arraystretch}{1.4}
\begin{tabular}{lll}
\toprule
\textbf{Option}
    & \textbf{Default}
    & \textbf{Accepted values}\\
\midrule
\texttt{color}
    & \texttt{black}
    & Valid xcolor expressions\\
\texttt{opacity}
    & \texttt{100}
    & Numbers from 0 to 100\\
\texttt{scale}
    & \texttt{1}
    & Positive numbers\\
\texttt{rotate}
    & \texttt{0}
    & Angles from 0 to 360\\
\bottomrule
\end{tabular}
\end{center}

Numerical values should be specified without units,
percentage signs, or degree symbols.

Use a decimal point rather than a decimal comma.

Numerical values are subject to the precision and
dimensional limits of PGF and TeX.

For example:

\begin{verbatim}
\klein[opacity=37.5, scale=2.5, rotate=22.5]
\end{verbatim}

\begin{center}
    \klein[opacity=37.5,scale=2.5,rotate=22.5]
\end{center}

% --------------------------------------------------
\section{Customization}
% --------------------------------------------------

\subsection{Color}

The \texttt{color} option changes the color of the
miniature, including its solid and dashed lines.

The default color is black.

Named colors and color expressions supported by
the \texttt{xcolor} package may be used.

\begin{center}
\renewcommand{\arraystretch}{2}
\begin{tabular}{lc}
\toprule
\textbf{Command} & \textbf{Result}\\
\midrule
\verb|\klein[color=black]|
    & \klein[color=black]\\
\verb|\klein[color=red]|
    & \klein[color=red]\\
\verb|\klein[color=blue]|
    & \klein[color=blue]\\
\verb|\klein[color=green]|
    & \klein[color=green]\\
\verb|\klein[color=orange]|
    & \klein[color=orange]\\
\verb|\klein[color=purple]|
    & \klein[color=purple]\\
\verb|\klein[color=cyan]|
    & \klein[color=cyan]\\
\verb|\klein[color=magenta]|
    & \klein[color=magenta]\\
\bottomrule
\end{tabular}
\end{center}

\subsubsection{Color Mixing}

The package also supports the color-mixing syntax
provided by \texttt{xcolor}.

For example, the expression \texttt{blue!50!red}
produces a mixture consisting of equal proportions
of blue and red.

\begin{center}
\renewcommand{\arraystretch}{2}
\begin{tabular}{lc}
\toprule
\textbf{Command} & \textbf{Result}\\
\midrule
\verb|\klein[color=blue!50!red]|
    & \klein[color=blue!50!red]\\
\verb|\klein[color=red!40!black]|
    & \klein[color=red!40!black]\\
\verb|\klein[color=green!50!black]|
    & \klein[color=green!50!black]\\
\verb|\klein[color=orange!70!black]|
    & \klein[color=orange!70!black]\\
\verb|\klein[color=blue!30!white]|
    & \klein[color=blue!30!white]\\
\bottomrule
\end{tabular}
\end{center}

\subsubsection{Custom Colors}

Users may define their own colors using the
\verb|\definecolor| command.

For example:

\begin{verbatim}
\definecolor{CustomGreen}{HTML}{2A9D8F}
\definecolor{CustomPurple}{HTML}{7952B3}

\klein[color=CustomGreen,scale=3]
\klein[color=CustomPurple,scale=3]
\end{verbatim}

\begin{center}
    \klein[color=CustomGreen,scale=3]
    \qquad
    \klein[color=CustomPurple,scale=3]
\end{center}

Any color model supported by \texttt{xcolor},
such as RGB, rgb, HTML, or CMYK, may be used
to define a named color.

Further information about color definitions
and color expressions can be found in the
\href{https://ctan.org/pkg/xcolor}
{\texttt{xcolor} documentation}.

% --------------------------------------------------
\subsection{Opacity}
% --------------------------------------------------

The \texttt{opacity} option controls the transparency
of the miniature.

Opacity is expressed as a percentage, with values
ranging from 0 to 100, inclusive:

\begin{itemize}
    \item \texttt{opacity=100}: fully opaque;
    \item \texttt{opacity=50}: partially transparent;
    \item \texttt{opacity=0}: fully transparent.
\end{itemize}

The default value is 100.

\begin{center}
\renewcommand{\arraystretch}{2}
\begin{tabular}{lc}
\toprule
\textbf{Command} & \textbf{Result}\\
\midrule
\verb|\klein[opacity=100]| & \klein[opacity=100]\\
\verb|\klein[opacity=80]|  & \klein[opacity=80]\\
\verb|\klein[opacity=60]|  & \klein[opacity=60]\\
\verb|\klein[opacity=40]|  & \klein[opacity=40]\\
\verb|\klein[opacity=20]|  & \klein[opacity=20]\\
\verb|\klein[opacity=0]|   & \klein[opacity=0]\\
\bottomrule
\end{tabular}
\end{center}

For a clearer comparison, the following miniatures
are displayed at a larger scale:

\begin{center}
\begin{tabular}{ccccc}
    \klein[opacity=100,scale=4] &
    \klein[opacity=75,scale=4]  &
    \klein[opacity=50,scale=4]  &
    \klein[opacity=25,scale=4]  &
    \klein[opacity=0,scale=4]   \\[0.3cm]
    100\% & 75\% & 50\% & 25\% & 0\%
\end{tabular}
\end{center}

Notice that a fully transparent miniature still
occupies space in the document.

Decimal values are also supported, for example:

\begin{verbatim}
\klein[opacity=37.5]
\end{verbatim}

Opacity is applied to the individual drawing paths.
Consequently, overlapping strokes may appear slightly
darker when partial opacity is used.

The opacity option should not be confused with color
mixing. A transparent color allows the background to
remain visible, whereas mixing a color with white
produces a lighter, potentially fully opaque color.

% --------------------------------------------------
\subsection{Scale}
% --------------------------------------------------

The \texttt{scale} option changes the size of the
miniature by multiplying its original dimensions
by a positive factor.

The default value is 1.

Values strictly between 0 and 1 reduce the miniature,
whereas values greater than 1 enlarge it.

Scaling is proportional: the dimensions, line
widths, and dash patterns are scaled together.

\begin{center}
\renewcommand{\arraystretch}{1.8}
\begin{tabular}{lc}
\toprule
\textbf{Command} & \textbf{Result}\\
\midrule
\verb|\klein[scale=0.25]| & \klein[scale=0.25]\\
\verb|\klein[scale=0.5]|  & \klein[scale=0.5]\\
\verb|\klein[scale=1]|    & \klein[scale=1]\\
\verb|\klein[scale=1.5]|  & \klein[scale=1.5]\\
\verb|\klein[scale=2]|    & \klein[scale=2]\\
\verb|\klein[scale=3]|    & \klein[scale=3]\\
\verb|\klein[scale=5]|    & \klein[scale=5]\\
\bottomrule
\end{tabular}
\end{center}

The following illustration compares several
different scaling factors:

\begin{center}
\begin{tabular}{cccccc}
    \klein[scale=0.5] &
    \klein[scale=1]   &
    \klein[scale=2]   &
    \klein[scale=3]   &
    \klein[scale=4]   &
    \klein[scale=5]   \\[0.3cm]
    $0.5$ & $1$ & $2$ & $3$ & $4$ & $5$
\end{tabular}
\end{center}

The scale must be strictly positive.
Zero and negative scaling factors are not accepted.

Extremely small or large scaling factors are subject
to the numerical precision of PGF and the dimensional
limits of the TeX engine.

In particular, a sufficiently small positive value
may be rounded to zero by the underlying numerical
calculation mechanism.

% --------------------------------------------------
\subsection{Rotation}
% --------------------------------------------------

The \texttt{rotate} option rotates the miniature
counterclockwise by a given angle, expressed in degrees.

The accepted values range from 0 to 360, inclusive,
with a default value of 0.

\begin{center}
\renewcommand{\arraystretch}{2}
\begin{tabular}{lc}
\toprule
\textbf{Command} & \textbf{Result}\\
\midrule
\verb|\klein[rotate=0]|   & \klein[rotate=0]\\
\verb|\klein[rotate=45]|  & \klein[rotate=45]\\
\verb|\klein[rotate=90]|  & \klein[rotate=90]\\
\verb|\klein[rotate=135]| & \klein[rotate=135]\\
\verb|\klein[rotate=180]| & \klein[rotate=180]\\
\verb|\klein[rotate=225]| & \klein[rotate=225]\\
\verb|\klein[rotate=270]| & \klein[rotate=270]\\
\verb|\klein[rotate=315]| & \klein[rotate=315]\\
\verb|\klein[rotate=360]| & \klein[rotate=360]\\
\bottomrule
\end{tabular}
\end{center}

The following illustration shows eight different
orientations of the miniature:

\begin{center}
\begin{tabular}{cccc}
    \klein[rotate=0,scale=4]   &
    \klein[rotate=45,scale=4]  &
    \klein[rotate=90,scale=4]  &
    \klein[rotate=135,scale=4] \\[0.3cm]
    $0^\circ$ & $45^\circ$ & $90^\circ$ & $135^\circ$
    \\[0.8cm]
    \klein[rotate=180,scale=4] &
    \klein[rotate=225,scale=4] &
    \klein[rotate=270,scale=4] &
    \klein[rotate=315,scale=4] \\[0.3cm]
    $180^\circ$ & $225^\circ$ & $270^\circ$ & $315^\circ$
\end{tabular}
\end{center}

Fractional angles are also supported:

\begin{verbatim}
\klein[rotate=22.5]
\end{verbatim}

A rotation of 360 degrees produces the same
orientation as a rotation of 0 degrees.

To obtain an orientation corresponding to a clockwise
rotation, use its equivalent counterclockwise angle.
For example, a clockwise rotation of 45 degrees
can be obtained using:

\begin{verbatim}
\klein[rotate=315]
\end{verbatim}

% --------------------------------------------------
\section{Combining Options}
% --------------------------------------------------

All customization options are independent and
may be combined in a single command.

For example:

\begin{verbatim}
\klein[color=blue, opacity=75, scale=4, rotate=30]
\end{verbatim}

\begin{center}
    \klein[color=blue,opacity=75,scale=4,rotate=30]
\end{center}

The options may be specified in any order.
Therefore, the following two commands are equivalent:

\begin{verbatim}
\klein[color=red, scale=3, rotate=45]
\klein[rotate=45, color=red, scale=3]
\end{verbatim}

\begin{center}
    \klein[color=red,scale=3,rotate=45]
    \qquad
    \klein[rotate=45,color=red,scale=3]
\end{center}

\subsection{Additional Examples}

\begin{verbatim}
\klein[color=purple, opacity=50, scale=4, rotate=45]
\end{verbatim}

\begin{center}
    \klein[color=purple,opacity=50,scale=4,rotate=45]
\end{center}

\begin{verbatim}
\klein[color=orange!70!black, scale=5, rotate=90]
\end{verbatim}

\begin{center}
    \klein[color=orange!70!black,scale=5,rotate=90]
\end{center}

\begin{verbatim}
\klein[color=CustomGreen, opacity=80, scale=4,
       rotate=315]
\end{verbatim}

\begin{center}
    \klein[color=CustomGreen,opacity=80,scale=4,rotate=315]
\end{center}

% --------------------------------------------------
\section{Scope of Customization}
% --------------------------------------------------

The customization options apply only to the
individual invocation of the \verb|\klein| command.

They do not modify the default appearance of
subsequent miniatures.

For example:

\begin{verbatim}
\klein[color=red, scale=3]
\klein
\klein[color=blue, rotate=90]
\klein
\end{verbatim}

\begin{center}
    \klein[color=red,scale=3]
    \qquad
    \klein
    \qquad
    \klein[color=blue,rotate=90]
    \qquad
    \klein
\end{center}

Each invocation starts with the default settings
before applying the options supplied by the user.

If the same option is specified more than once
within a single invocation, the last assignment
takes precedence.

For example:

\begin{verbatim}
\klein[color=red, color=blue]
\end{verbatim}

\begin{center}
    \klein[color=red,color=blue]
\end{center}

The resulting miniature is blue.

% --------------------------------------------------
\section{Using the Miniature in Documents}
% --------------------------------------------------

\subsection{Inline Usage}

The default miniature is designed to be inserted
directly into ordinary text.

For example:

\begin{verbatim}
The Klein bottle \klein\ is a non-orientable surface.
\end{verbatim}

\begin{quote}
The Klein bottle \klein\ is a non-orientable surface.
\end{quote}

The miniature may also be customized within a
paragraph:

\begin{verbatim}
The Klein bottle \klein[color=blue]
is a classical object in topology.
\end{verbatim}

\begin{quote}
The Klein bottle \klein[color=blue]
is a classical object in topology.
\end{quote}

When the command is used without an optional
argument, an explicit control space, written as
\verb|\ |, may be used to preserve the following
interword space.

The command produces a graphical illustration,
rather than a character from a mathematical font
or a native mathematical symbol.

Its size is determined by the \texttt{scale}
option and is not automatically adjusted to
mathematical styles such as script style or
scriptscript style.

Large scaling factors and rotations may affect
the height and depth of the surrounding text line.

\subsection{Displayed Illustrations}

Larger scaling factors can be used to produce
standalone illustrations.

\begin{verbatim}
\begin{center}
    \klein[color=blue!70!black, scale=8]
\end{center}
\end{verbatim}

\begin{center}
    \klein[color=blue!70!black,scale=8]
\end{center}

Since the miniature is drawn using TikZ,
it is not based on an external bitmap image.
In PDF output, it retains its vector-based
appearance when enlarged.

\subsection{Figures and Captions}

The miniature may also be placed inside a
standard figure environment.

\begin{verbatim}
\begin{figure}[ht]
    \centering
    \klein[color=ManualBlue, scale=5]
    \caption{A representation of the Klein bottle.}
\end{figure}
\end{verbatim}

\begin{center}
    \klein[color=ManualBlue,scale=5]
    \par\vspace{0.2cm}
    {\small Figure example: A representation
    of the Klein bottle.}
\end{center}

In this example, \texttt{ManualBlue} is a custom
color defined in the preamble of this documentation.

\subsection{Use in Tables}

The miniature can also be used inside table cells,
as demonstrated throughout this documentation.

For example:

\begin{verbatim}
\begin{tabular}{lc}
    Default miniature & \klein\\
    Blue miniature    & \klein[color=blue]\\
    Rotated miniature & \klein[rotate=90]
\end{tabular}
\end{verbatim}

\begin{center}
\renewcommand{\arraystretch}{1.5}
\begin{tabular}{lc}
    Default miniature & \klein\\
    Blue miniature    & \klein[color=blue]\\
    Rotated miniature & \klein[rotate=90]
\end{tabular}
\end{center}

\subsection{Use with Beamer}

The package can also be used in presentations
created with the
\href{https://ctan.org/pkg/beamer}
{\texttt{beamer}} document class.

For example:

\begin{verbatim}
\documentclass{beamer}
\usepackage{kleinbottle}

\begin{document}

\begin{frame}{The Klein Bottle}

    The Klein bottle is a non-orientable surface.

    \begin{center}
        \klein[color=blue, scale=5]
    \end{center}

\end{frame}

\end{document}
\end{verbatim}

The same customization options are available
in Beamer presentations.

\subsection{Use in PDF Bookmarks}

The command generates a graphical object and
therefore cannot be represented directly as
plain text in PDF bookmarks.

When using the miniature in a section title
that generates a PDF bookmark, users of
\texttt{hyperref} may provide an alternative
text representation using \verb|\texorpdfstring|.

For example:

\begin{verbatim}
\section{\texorpdfstring{\klein\ The Klein Bottle}
                       {The Klein Bottle}}
\end{verbatim}

% --------------------------------------------------
\section{Error Handling}
% --------------------------------------------------

The package checks the numerical values supplied
to the \texttt{opacity}, \texttt{scale}, and
\texttt{rotate} options.

Values outside the accepted ranges generate
package errors.

The following examples illustrate invalid
commands and the reasons why they are rejected.

\begin{center}
\renewcommand{\arraystretch}{1.5}
\begin{tabular}{ll}
\toprule
\textbf{Invalid command} & \textbf{Reason}\\
\midrule
\verb|\klein[opacity=-10]| & Negative opacity\\
\verb|\klein[opacity=120]| & Opacity above 100\\
\verb|\klein[scale=0]|     & Zero scaling factor\\
\verb|\klein[scale=-2]|    & Negative scaling factor\\
\verb|\klein[rotate=-45]|  & Negative rotation\\
\verb|\klein[rotate=400]|  & Rotation above 360 degrees\\
\bottomrule
\end{tabular}
\end{center}

If compilation continues after a numerical
range error, the corresponding option is
replaced by its default value.

Undefined colors and unrecognized option names
are handled by the underlying \texttt{xcolor}
and \texttt{pgfkeys} mechanisms.

Non-numeric or empty numerical arguments may
produce errors from PGFMath. The fallback
mechanism described above applies specifically
to numerical range errors.

Users should correct invalid arguments rather
than relying on compilation to continue after
an error.

% --------------------------------------------------
\section{Minimal Working Example}
% --------------------------------------------------

The following is a complete example of a document
using the \texttt{kleinbottle} package, together with
its corresponding output.

\medskip

\noindent
\begin{minipage}[t]{0.58\textwidth}

    \vspace{0pt}

    \textbf{Source code}

    \small

\begin{verbatim}
\documentclass{article}

\usepackage{kleinbottle}

\definecolor{MyColor}{HTML}{234A76}

\begin{document}

The Klein bottle \klein\ is a
compact non-orientable surface.

\begin{center}

    \klein[color=MyColor, scale=5]

    \klein[color=red, opacity=50,
           scale=3, rotate=45]

\end{center}

\end{document}
\end{verbatim}

\end{minipage}%
\hfill
\begin{minipage}[t]{0.38\textwidth}

    \vspace{0pt}

    \textbf{Output}

    \medskip

    \noindent
    \fbox{%
        \begin{minipage}{
            \dimexpr\linewidth-2\fboxsep-2\fboxrule\relax}

            \definecolor{MyColor}{HTML}{234A76}

            \small
            \vspace{0.3cm}

            The Klein bottle \klein\ is a
            compact non-orientable surface.

            \begin{center}

                \klein[color=MyColor, scale=5]

                \klein[color=red, opacity=50,
                       scale=3, rotate=45]

            \end{center}

            \vspace{0.3cm}

        \end{minipage}%
    }

\end{minipage}

\medskip

The output shown above reproduces the contents
of the example document within a smaller
typesetting area.

% --------------------------------------------------
\section{License and Maintenance}
% --------------------------------------------------

Copyright \textcopyright{} 2026 Benjamim Lima Neto.

This work may be distributed and/or modified under
the conditions of the \LaTeX{} Project Public License
(LPPL), either version 1.3c of this license or,
at your option, any later version.

The latest version of this license is available at:

\begin{center}
    \url{https://www.latex-project.org/lppl.txt}
\end{center}

A copy of the license is included in the accompanying
\texttt{LICENSE} file.

This work has the LPPL maintenance status
\texttt{maintained}.

The author and current maintainer of the package
is Benjamim Lima Neto.

Bug reports and suggestions may be submitted using
the contact information provided in the accompanying
\texttt{README.md} file.

\subsection{Files Included in the Distribution}

The package distribution contains the following files:

\begin{center}
\renewcommand{\arraystretch}{1.4}
\begin{tabular}{ll}
\toprule
\textbf{File} & \textbf{Description}\\
\midrule
\texttt{kleinbottle.sty} & Package source code\\
\texttt{kleinbottle.tex} & Documentation source code\\
\texttt{kleinbottle.pdf} & Compiled documentation\\
\texttt{README.md}       & Package information\\
\texttt{LICENSE}         & LPPL license text\\
\bottomrule
\end{tabular}
\end{center}

The package source, documentation source, README file,
and derived documentation constitute the Work under
the LPPL.

The accompanying license text is distributed verbatim
and retains its original copyright notice.

% --------------------------------------------------
\section{Version History}
% --------------------------------------------------

\noindent
\textbf{Version 1.0 -- November 26, 2026}

\begin{itemize}

    \item Initial release.

    \item Introduction of the \verb|\klein| command.

    \item Support for named colors, color mixing,
          and custom color definitions.

    \item Percentage-based opacity customization.

    \item Proportional scaling of the illustration,
          including line widths and dash patterns.

    \item Counterclockwise rotation in degrees.

    \item Independent customization of individual
          miniatures.

    \item Validation of numerical option ranges.

    \item User documentation with illustrated examples.

\end{itemize}

\vfill

\begin{center}
    \klein[color=ManualBlue,scale=2]
    \qquad
    \klein[color=CustomGreen,scale=2]
    \qquad
    \klein[color=CustomPurple,scale=2]
\end{center}

\end{document}
